\unnumberedchapter{MỞ ĐẦU}
\label{chap:introduction}

Biển báo giao thông cung cấp thông tin về quy định, cảnh báo và hướng đi trên đường. Nhận diện biển báo từ ảnh hoặc video gồm tìm vị trí biển, xác định loại biển và đọc nội dung chữ, số khi có. Bài toán gặp nhiều khó khăn: biển ở xa rất nhỏ, ánh sáng yếu hoặc ảnh mờ làm thay đổi màu sắc và hình dạng, nhiều loại biển chỉ khác nhau ở ký hiệu bên trong, còn nền đường phố có nhiều vật thể cùng màu với biển. Các công trình về biển báo Việt Nam gần đây chủ yếu dùng họ YOLO trên dữ liệu riêng~\cite{tung2024vietnam,cu2024vietnamese}, ít so sánh trực tiếp với phương pháp xử lý ảnh truyền thống.

Từ thực tiễn trên, nhóm chọn đề tài ``Hệ thống nhận diện biển báo giao thông''. Hệ thống kết hợp học sâu (YOLO11 phát hiện và phân loại, VietOCR đọc chữ và số) với thị giác máy tính truyền thống (đánh giá chất lượng và tăng cường ảnh, theo vết bằng bộ lọc Kalman, đặc trưng màu, hình dạng, HOG kết hợp SVM). Đề tài có bốn mục tiêu: (M1) xây dựng hệ thống hoàn chỉnh nhận ảnh, video, webcam và trả về vị trí, tên lớp, giá trị số, ý nghĩa theo QCVN cho 27 lớp biển; (M2) đánh giá riêng từng phương pháp và so sánh trực tiếp đặc trưng thủ công với học sâu; (M3) cải thiện kết quả với biển nhỏ và ảnh tối mà không bổ sung dữ liệu; (M4) đánh giá đáng tin cậy, có khoảng tin cậy và tái lập được.

Phạm vi đề tài là 27 lớp trong bộ dữ liệu biển báo Việt Nam đã chọn, giữ cố định trong suốt quá trình thực hiện; không xử lý đèn tín hiệu, vạch kẻ đường hay điều khiển phương tiện. Báo cáo gồm ba chương: Chương~\ref{chap:cosolythuyet} trình bày cơ sở lý thuyết và các nghiên cứu liên quan; Chương~\ref{chap:chuong2} trình bày dữ liệu, kiến trúc hệ thống và phương pháp; Chương~\ref{chap:chuong3} trình bày kết quả thực nghiệm, chương trình demo và thảo luận.
